\documentclass[11pt,letterpaper]{article}
\usepackage[margin=1in]{geometry}\usepackage{amsmath,amssymb,booktabs,array,microtype}\usepackage[T1]{fontenc}\usepackage{lmodern}
\usepackage[colorlinks=true,linkcolor=blue,citecolor=blue,urlcolor=blue]{hyperref}
\setlength{\parindent}{0pt}\setlength{\parskip}{5pt}\sloppy
\newcolumntype{L}[1]{>{\raggedright\arraybackslash}p{#1}}
\title{Observing Case: Does the Galactic Acceleration Scale Evolve?\\ Deep-Regime Rotators at $z\simeq2$--2.6 \\ \large Version 2: corrected target ledger, separated rivals, honest odds}
\author{Carl P.\ Zimmerman\\ \small Briar Creek Tech\\ \small AI-assisted research programme; not peer reviewed}
\date{11 September 2026; version 2, 25 September 2026}
\begin{document}\maketitle
\begin{abstract}
Companion to the pre-registered measurement (doi:10.5281/zenodo.22563139), written for observers. This version
supersedes version 1, whose candidate ledger carried values that contradict their own cited sources, and which
understated what a single galaxy can decide. Three readings of the radial-acceleration scale $a_0$ make
different zero-parameter predictions for its value at redshift $z$, quoted separately here as
$\Delta\equiv\log_{10}[a_0(z)/a_0(0)]$ at $z=2.5$: constant ($0.00$), the emergent scale of $\Lambda$CDM haloes
($+0.23$ to $+0.46$ dex over halo masses $10^{11}$--$10^{13}\,M_\odot$ and two concentration--mass relations), and
$a_0\propto H(z)$ ($+0.58$ dex). With intrinsic and halo-to-halo scatter included, one galaxy measured to
$\pm0.13$ dex gives odds of about 4:1 between constant and $\Lambda$CDM-emergent, not 20:1. That decision needs
three galaxies at $\pm0.10$ dex, or four at $\pm0.20$. After every row of the ledger was re-read from its primary
paper, no known galaxy passes the gate. One rotation-supported, low-acceleration galaxy at $z\geq2$ is a viable
first target: OLAS M0717-02064 at $z=2.07$. Its published kinematics favour none of the three readings. For
velocity data alone, the per-object precision is gas-limited at $\approx0.43$ dex. A gas mass good to $\pm0.12$
dex brings it to $\approx0.26$ dex.
\end{abstract}

\section*{What changed from version 1}
\begin{enumerate}\itemsep2pt
\item \textbf{The ledger was wrong.} Several rows contradicted their cited sources. The replacement,
\texttt{highz\_target\_ledger\_verified\_2026.py}, re-reads every value from its primary paper. The largest error:
A68-HLS115 lies at $z=1.5869$ with $M_\star=8.1\times10^9\,M_\odot$ (Dessauges-Zavadsky et al.\ 2015). It is not at
$z=2.49$, and it is not a deep-regime candidate.
\item \textbf{Two rivals were merged.} Version 1 called the $\Lambda$CDM-emergent rise ``$\propto H(z)$''. The
emergent law, $E(z)^{4/3}c^2/f(c)$, gives $+0.33$ dex at $z=2.5$ for $10^{12}\,M_\odot$ haloes with Dutton--Macci\`o
concentrations. The $H(z)$ law gives $+0.58$ dex. They are now separate hypotheses.
\item \textbf{One galaxy does not decide.} The ``20:1 with one galaxy at $\pm0.13$ dex'' ignored the relation's
intrinsic scatter and the halo-to-halo concentration scatter. The corrected odds are in Sec.~3.
\item \textbf{Sign convention.} Version 1 defined $\Delta_{\rm BTFR}=\log M_b-4\log V_f-C_0$ but quoted a rise as
positive. For a rising $a_0$ that quantity is negative. This version uses $\Delta=\log_{10}[a_0(z)/a_0(0)]$
throughout.
\item \textbf{A caveat on the framework's own prediction.} The constant law is the framework's, but its observable
zero point is not yet sharp. A retained cold halo around $z\simeq2.5$ galaxies (internal note L189) raises the
apparent, baryons-only zero point by $\approx+0.2$ dex in the one example worked there, a $3\times10^{11}\,M_\odot$
halo around a rotator with $10^{10}\,M_\odot$ of baryons. That shift points toward the lower end of the
$\Lambda$CDM range. The correction has not been computed. Until it is, this case tests the three hypotheses above
and not the framework as a whole.
\end{enumerate}

\section{The three hypotheses}
\begin{center}\begin{tabular}{L{5.2cm}ccc}\toprule
$\Delta=\log_{10}[a_0(z)/a_0(0)]$ & constant & $\Lambda$CDM emergent & $a_0\propto H(z)$\\\midrule
$z=2.07$ (target 1) & $0.00$ & $+0.17$ to $+0.26$ & $+0.50$\\
$z=2.5$ & $0.00$ & $+0.23$ to $+0.46$ & $+0.58$\\
$z=2.62$ (target 2) & $0.00$ & $+0.24$ to $+0.35$ & $+0.60$\\\bottomrule\end{tabular}\end{center}
At target redshifts the $\Lambda$CDM ranges are for $10^{11}$--$10^{12}\,M_\odot$ haloes (DM14). At $z=2.5$ the range
spans $10^{11}$--$10^{13}\,M_\odot$ and both the DM14 and D08 concentration--mass relations. In the deep regime,
$V_c\propto a_0^{1/4}$ at fixed baryonic mass. The local calibration is frozen before any high-redshift value is
examined, and both local footings ($9.36\times10^{-11}$ and $1.13\times10^{-10}\,{\rm m\,s^{-2}}$) are always reported.

\section{The gate (unchanged except the window)}
A galaxy enters the decisive sample only if all of the following hold:
\begin{enumerate}\itemsep0pt
\item $g_{\rm bar}(R_{\rm out})<0.3\,a_0$ at the upper end of its mass range;
\item $V_{\rm rot}/\sigma>1.5$ in a forward-modelled source-plane fit, with inclination, dispersion, beam smearing
and asymmetric drift fitted together;
\item $\delta\log\mu\lesssim0.05$ dex (tightened from 0.08: the magnification enters the inverted $a_0$ with weight
$A_{\rm obs}/2\approx1.3$);
\item $35^\circ<i<75^\circ$, with at least 5 source-plane resolution elements along the major axis.
\end{enumerate}
The redshift window is relaxed from $2.32<z<3.12$ to $z\geq2$: at $z=2.07$, H$\alpha$ falls in NIRSpec G235H,
[O\,III] and H$\beta$ fall in G140H, and CO(3--2) falls at 112.6 GHz, inside ALMA Band 3.

\section{What one, two, three galaxies decide}
Near the gate the inversion $a_0=g_{\rm bar}/y$ of the relation $\nu(y)=1/(1-e^{-\sqrt y})$ amplifies errors. At
$y=0.2$, a mass error is amplified $\times1.52$ and an error in $g_{\rm obs}$ (twice the log-velocity error)
$\times2.52$. The deep-limit estimator $4\log V-\log M_b$ understates both. Two scatters add to the measurement
error: the relation's intrinsic scatter (0.034 dex in $g_{\rm obs}$, i.e.\ 0.08--0.09 dex in $a_0$), and, on the
$\Lambda$CDM side only, $\pm0.13$ dex in each object's expectation from halo-to-halo concentration scatter.
Expected odds between constant and $\Lambda$CDM-emergent ($\Delta=0.33$) at $z=2.5$:
\begin{center}\begin{tabular}{lcccc}\toprule
$\sigma_{\rm meas}$ per galaxy & $N=1$ & $N=2$ & $N=3$ & $N=4$\\\midrule
0.10 dex & 5.0:1 & 24.6:1 & 122:1 & 606:1\\
0.13 dex & 3.8:1 & 14.1:1 & 53:1 & 199:1\\
0.20 dex & 2.3:1 & 5.4:1 & 12.5:1 & 28.9:1\\\bottomrule\end{tabular}\end{center}
\textbf{Velocity data alone} (gas from the resolved star-formation law, $\pm0.30$ dex; $V_c$ to 5\%; inclination
$\pm8^\circ$): $\sigma(\log a_0)\approx0.43$ dex per object at $y=0.2$. Three such objects separate constant from
$a_0\propto H(z)$ at $\approx2\sigma$ and do not separate constant from $\Lambda$CDM-emergent ($\approx0.7\sigma$).
\textbf{With a gas mass good to $\pm0.12$ dex} (CO or dust continuum, with metallicity-dependent conversion), the
error is $\approx0.26$ dex per object. Three objects then reach $\approx3.5\sigma$ on constant versus $H(z)$.
Constant versus $\Lambda$CDM-emergent still needs the tighter budgets in the table. \emph{The gas mass is the hard
half of the measurement.}

\section{Candidates (verified ledger)}
\textbf{OLAS M0717-02064} ($z=2.07$, $\mu=6.48$, RA/Dec 07:17:39.125 +37:44:18.45; Hirtenstein et al.\ 2019,
Tables 1--2) is lens-corrected $\log M_\star=8.08\pm0.06$, with $\Delta v=142$ km s$^{-1}$,
$\sigma_{\rm local}=42.9$ km s$^{-1}$, $v/\sigma=1.64$ (class ``rotation'') and H$\alpha$ SFR 4.5
$M_\odot\,{\rm yr^{-1}}$. It is the only $z\geq2$ object that is both rotation-supported and plausibly
deep-regime, and it has no JWST spectroscopy yet (the MAST archive, 25 September 2026, shows NIRCam imaging only).
\emph{Face value.} Invert the full relation at $R=1.5$--3.5 kpc on both footings. A constant scale needs
$M_{\rm gas}/M_\star\approx6$--36 (depletion time 0.17--0.95 Gyr), and $a_0\propto H(z)$ needs $\approx3$--17.
Both are plausible, so \textbf{the published data favour none of the hypotheses}. (An intermediate calculation
with the deep-limit formula suggested a lean against constancy. It overstated the required gas by
$\times1.5$--3.4 and is withdrawn.) There is also a risk: if the disc is inclined and the scale is constant, the
galaxy sits at $y\approx0.4$--1.8 within 3.5 kpc and could fail gate 1 unless the velocity field reaches further
out.

\textbf{Abell 68 C4} ($z=2.622$, $\mu\approx35$--46) has two published stellar masses 0.7 dex apart
($\log M_\star=8.7$, Richard et al.\ 2007; $8.0$, Richard et al.\ 2011). It is deep-regime-plausible only at the
lower value. No kinematics exist. \textbf{Abell 68 C20b} ($\mu\approx158$, range 83--302) fails the lens gate.
\textbf{SL2S J0217} ($z=1.844$, $\mu=17.3$) has $\log M_\star=8.26$ and a metallicity of $0.05\,Z_\odot$ (Berg et al.\
2018). Existing ALMA [C\,II] (tentative) and PdBI CO(2--1) (undetected) data constrain its cold gas (Rybak et al.\
2021), but its CO is expected to be dark. \textbf{Controls:} A1689B11 ($z=2.54$; H$\alpha$ detected only to
$\approx1.7$ kpc; $V/\sigma=9$--13 in Yuan et al.\ 2017) and zC-400569 ($z=2.24$; $V_{\rm rot}=254\pm41$ km
s$^{-1}$, $V/\sigma\gtrsim17$), both high-acceleration. They must return the local zero point, where all three
hypotheses agree.

\section{Decision rule}
The rule is stated per hypothesis and is unchanged in form. Measured $\Delta$ near $0.00$ keeps the constant law
and excludes the others at the achieved odds. Near the $\Lambda$CDM range, constancy is disfavoured. Near
$+0.5$--$0.6$, both constancy and the emergent law are disfavoured. A value inconsistent with all three is reported
as such. Odds are quoted with the scatters of Sec.~3, never from the measurement error alone. A null on the gate
(no galaxy qualifies) is itself informative about the deep-regime population at $z\geq2$.

\medskip\small\textbf{Provenance.} doi:10.5281/zenodo.22563139 (pre-registration). The repository scripts are
\texttt{prep\_2026/a0z\_crossscale/highz\_target\_ledger\_verified\_2026.py} (verified ledger) and
\texttt{prep\_2026/a0z\_crossscale/a0z\_deepmond\_ifu\_forecast\_2026.py} (hypotheses, face value, per-object
precision; 14/14 checks). The odds table uses the same scatter model as the programme's step-by-step
manuscript (\texttt{paper\_numbers.py}, S4).
\end{document}
